\documentclass[11pt,a4paper]{article}
\usepackage[T1]{fontenc}
\usepackage[utf8]{inputenc}
\usepackage[ngerman]{babel}
\usepackage[margin=2.5cm]{geometry}
\usepackage[hidelinks]{hyperref}
\usepackage{array}
\setlength{\parindent}{0pt}
\setlength{\emergencystretch}{2em}
\renewcommand{\arraystretch}{1.25}
\sloppy
\begin{document}
\shorthandoff{"}
\begin{center}
  {\Large\textbf{Webtechnologien -- Übungsblatt 5}}
\end{center}
\section*{Aufgabe 3: Promises II}
\textbf{Aufgabe:} Erklären, was \texttt{5\_3.js} macht und warum die
Ausgabe in dieser Reihenfolge erscheint. Danach
5+3+7+14+23+25+1000+1 mit \texttt{delayedCalc(...)} schrittweise in genau
dieser Reihenfolge berechnen, einmal mit einer Kette aus
\texttt{.then(...)} und \texttt{.catch(...)} und einmal mit
\texttt{async}/\texttt{await} und \texttt{try}/\texttt{catch}.
\texttt{delayedCalc} selbst darf nicht geändert werden.

\medskip
\textbf{a) Was macht der gegebene Code?}
\begin{itemize}
  \item \texttt{delayedCalc(a, b)} gibt \emph{sofort} ein Promise zurück,
    also ein Versprechen, dass das Ergebnis später kommt.
  \item Ist \texttt{a > b}, startet ein \texttt{setTimeout} mit einer
    zufälligen Wartezeit von 0 bis 1999\,ms. Danach wird
    \glqq a + b = c\grqq{} ins Event Log geschrieben und das Promise mit
    \texttt{resolve(c)} erfüllt.
  \item Ist \texttt{a <= b}, wird das Promise \emph{sofort} (ohne Wartezeit)
    mit \texttt{reject('error: b larger than a')} abgelehnt und
    \glqq rejected\grqq{} geschrieben.
  \item Unten wird \texttt{delayedCalc} viermal direkt hintereinander
    aufgerufen, ohne auf ein Ergebnis zu warten. Alle vier laufen also
    gleichzeitig, und keine Rechnung bekommt das Ergebnis der vorigen.
\end{itemize}
\textbf{Warum diese Reihenfolge?} \texttt{delayedCalc(7,14)} und
\texttt{delayedCalc(23,25)} lehnen sofort ab, darum stehen zweimal
\glqq rejected\grqq{} immer ganz oben. In der Konsole erscheint zweimal
\texttt{Uncaught (in promise) error: b larger than a}, weil niemand die
Ablehnung mit \texttt{catch} abfängt. \texttt{5 + 3 = 8} und
\texttt{1000 + 1 = 1001} warten zufällig lange; wer die kürzere Wartezeit
hat, steht zuerst. Deshalb ändert sich die Reihenfolge dieser beiden
Zeilen beim Neuladen.

\medskip
\textbf{Ergebnis} (\texttt{ex3.html} unverändert): \texttt{5\_3.js} ist die
Originaldatei. Die vier alten Aufrufe sind nur auskommentiert und mit
\texttt{[3b ENTFERNT]} markiert; alles Neue steht nach fünf Leerzeilen am
Dateiende zwischen \texttt{[NEU]} und \texttt{[ENDE NEU]}.
\begin{itemize}
  \item \textbf{b)} Kette: \texttt{delayedCalc(5, 3)}, dann
    \texttt{.then(result => delayedCalc(result, 7))} usw. Jedes
    \texttt{.then} gibt das neue Promise zurück, darum wartet die Kette,
    bis die Rechnung fertig ist, und \texttt{result} ist immer das
    Ergebnis davor. Bei \texttt{77 + 1000} ist \texttt{b} größer, das
    Promise wird abgelehnt, alle weiteren \texttt{.then} werden
    übersprungen und das eine \texttt{.catch(handleError)} am Ende fängt
    den Fehler. Diese Version steht als Kommentar in der Datei, weil nur
    eine Version gleichzeitig laufen darf (sonst mischen sich die
    Ausgaben).
  \item \textbf{c)} \texttt{async function calcWithAwait()}: Jede Zeile
    \texttt{result = await delayedCalc(result, ...)} wartet, bis das
    Promise fertig ist, und nimmt dann den Wert. Ein abgelehntes Promise
    wirkt bei \texttt{await} wie ein geworfener Fehler und springt in den
    \texttt{catch}-Block, der \texttt{handleError(error)} aufruft. Diese
    Version ist aktiv.
  \item Umschalten auf b): nur bei den 8 Code-Zeilen von b) (von
    \texttt{delayedCalc(5, 3)} bis \texttt{.catch(handleError);}) das
    \texttt{//} am Zeilenanfang entfernen und den Aufruf \texttt{calcWithAwait();} auskommentieren.
\end{itemize}

\textbf{Hinweis:} \texttt{1000 + 1} wird nie gerechnet, weil die Kette
schon bei \texttt{77 + 1000} abbricht. Das entspricht der verlangten
Ausgabe in der Angabe, die mit \glqq rejected\grqq{} endet.
\texttt{handleError} zeigt den Fehlertext als \texttt{alert}.

\medskip
\textbf{Test:} In Chrome, \texttt{ex3.html} direkt geöffnet:
\begin{center}
\begin{tabular}{|>{\raggedright\arraybackslash}p{4.4cm}|>{\raggedright\arraybackslash}p{9.3cm}|}
\hline
\textbf{Version} & \textbf{Ergebnis} \\ \hline
Original (a) & zuerst 2x rejected, dann 5 + 3 = 8 und 1000 + 1 = 1001 in
  wechselnder Reihenfolge, 2x Uncaught (in promise) \\ \hline
b) \texttt{.then}-Kette & 5 + 3 = 8, 8 + 7 = 15, 15 + 14 = 29,
  29 + 23 = 52, 52 + 25 = 77, rejected; ein alert
  \glqq error: b larger than a\grqq \\ \hline
c) \texttt{async}/\texttt{await} & genau dieselbe Ausgabe und derselbe
  alert \\ \hline
\end{tabular}
\end{center}
Bei b) und c) keine Fehler in der Konsole; mehrmals neu geladen, die
Reihenfolge bleibt immer gleich.

\medskip
{\small \textbf{Hilfsmittel:} Claude Code (Claude Opus 5.5) für Erstellung, Prüfung und Dokumentation.}
\end{document}
